%% shortname: Equivalent and Compensating Variation
\documentclass[11pt]{article}
\usepackage{microstyle}

\begin{document}

\lessonheader{10}{Chapter 3 (continued): welfare evaluation --- equivalent variation, compensating variation and consumer surplus}{}


\section{The money-metric welfare change}

Revisiting the money-metric indirect utility function to derive the welfare impact
of a price change from $p^0$ to $p'$, given the benchmark price vector $\bar p$:
\[ v_m(p',w\mid\bar p)-v_m(p^0,w\mid\bar p)
   =e\bigl(\bar p,v(p',w)\bigr)-e\bigl(\bar p,v(p^0,w)\bigr). \]
If positive, then the consumer is better off under $p'$ than under $p^0$.

The \emph{sign} is the same regardless of $\bar p$, but the \emph{size} does depend
on $\bar p$, simply because we are using different prices to evaluate the welfare
impact.

Two benchmarks are particularly natural: $\bar p=p^0$ and $\bar p=p'$, i.e.\
evaluating the change at the old prices or at the new prices.

Throughout, write $u^0=v(p^0,w)$ and $u'=v(p',w)$ for the utility levels before
and after the price change. By duality,
$e(p^0,u^0)=w=e(p',u')$.

\section{Equivalent and compensating variation}

\topic{Equivalent variation ($\bar p=p^0$)}
If $\bar p=p^0$, the equivalent variation is
\begin{align*}
  EV(p^0,p',w)&=e\bigl(p^0,v(p',w)\bigr)-e\bigl(p^0,v(p^0,w)\bigr)\\
              &=e(p^0,u')-e(p^0,u^0)\\
              &=e(p^0,u')-w .
\end{align*}
If, under $p^0$, the consumer receives $e(p^0,u')-w$, then he can afford $u'$,
because his total wealth is $e(p^0,u')$. So this is the amount the consumer would
have to be paid at the old prices $p^0$ to be willing to forgo the price change. Hence it is the amount
that is \emph{equivalent} to the price change.

\topic{Compensating variation ($\bar p=p'$)}
If $\bar p=p'$, the compensating variation is
\begin{align*}
  CV(p^0,p',w)&=e\bigl(p',v(p',w)\bigr)-e\bigl(p',v(p^0,w)\bigr)\\
              &=e(p',u')-e(p',u^0)\\
              &=w-e(p',u^0).
\end{align*}
If, at the new prices $p'$, you give up $w-e(p',u^0)$, you end up with wealth
$e(p',u^0)$, which is just enough to get you back to the old level of utility $u^0$.
It is the amount needed to \emph{compensate} for the change in prices.

\runin{Example.} $p_i'<p_i^0$, all else equal. The agent is better off: $u'>u^0$.
Both $EV$ and $CV$ have to be positive, but they are likely to be different.


\section{Comparing EV and CV}

Let's try to compare the size of $EV$ and $CV$. From Shephard's lemma, there is a
link between the expenditure function and Hicksian demand. Let us use that link,
together with the link between Hicksian and Marshallian demand in the Slutsky
equation.

Recall
\begin{align*}
  EV(p^0,p',w)&=e(p^0,u')-w\\
              &=e(p^0,u')-e(p',u').
\end{align*}
If only $p_i$ changes, then using the fundamental theorem of calculus,
\[ EV(p^0,p',w)=\int_{p_i'}^{p_i^0}\pd{e(p,u')}{p_i}\,dp_i
   =\int_{p_i'}^{p_i^0}h_i\bigl(p_1^0,p_2^0,\dots,p_i,\dots,p_L^0,\,u'\bigr)\,dp_i . \]
\begin{flatlist}
  \item[$\Rightarrow$] $EV$ = area under the Hicksian demand curve evaluated at $u'$;
  \item[] $CV$ = area under the Hicksian demand curve evaluated at $u^0$.
\end{flatlist}

We need to compare Hicksian demand at different utility levels. We take a few steps.

\topic{Step 1. Hicksian vs.\ Marshallian demand for fixed $u$}
The Slutsky equation, evaluated at $u=v(p,w)$ (equivalently $w=e(p,u)$):
\[ \pd{h_i(p,u)}{p_i}=\pd{x_i(p,w)}{p_i}+\pd{x_i(p,w)}{w}\,x_i(p,w)\le 0 , \]
because of concavity of the expenditure function: the negative own-price substitution
effect --- Hicksian demand slopes down.

At this point, add the assumption that good $i$ is a \emph{normal} good:
\[ \pd{x_i(p,w)}{w}>0
   \quad\Longrightarrow\quad \pd{x_i(p,w)}{p_i}<0 , \]
so Marshallian demand slopes down: it is an \emph{ordinary} good. Moreover, at the same point
($u=v(p,w)$) and whenever $x_i(p,w)>0$,
\[ \pd{h_i(p,u)}{p_i}>\pd{x_i(p,w)}{p_i}. \]
The left-hand side is less negative than the right-hand side.
\begin{flatlist}
  \item[$\Rightarrow$] Hicksian demand is less responsive to its own price than
        Marshallian demand;
  \item[$\Rightarrow$] in the usual graph, with price on the vertical axis (the inverse
        demand curve), the Hicksian curve is steeper.
\end{flatlist}

\topic{Step 2. Where the curves cross}
An implication is that Hicksian and Marshallian demand can cross only once (with the other prices fixed at
$p_{-i}^0$ and $w$ fixed). What can
we say about the price at which the crossing occurs, and how does it depend on
utility?

For concreteness, assume $p_i'<p_i^0$, and thus $u'>u^0$. From duality,
\begin{align*}
  h_i(p^0,u^0)&=x_i\bigl(p^0,e(p^0,u^0)\bigr)=x_i(p^0,w),\\
  h_i(p',u') &=x_i\bigl(p',e(p',u')\bigr)=x_i(p',w)>x_i(p^0,w),
\end{align*}
since the good is an ordinary good. So $h_i(\cdot,u^0)$ crosses Marshallian demand
at $p_i^0$, and $h_i(\cdot,u')$ crosses it at $p_i'$.

\begin{center}
\begin{tikzpicture}[scale=1.05]
  \ecaxes{6.2}{4.9}
  \node[ecnode,above] at (0,4.9) {price};
  \node[ecnode,right] at (6.2,0) {quantity};
  % Marshallian demand p = 4/x
  \draw[budget] plot[domain=0.95:5.6,samples=60] (\x,{4/\x});
  \node[ecnode,right] at (5.6,0.71) {$x_i(\cdot,w)$};
  % Hicksian demand at u^0: p = 4.8/x^2
  \draw[budget,dashed] plot[domain=1.05:2.6,samples=60] (\x,{4.8/(\x*\x)});
  \node[ecnode,anchor=north west] at (2.56,0.6) {$h_i(\cdot,u^0)$};
  % Hicksian demand at u': p = 10/x^2
  \draw[budget,dashed] plot[domain=1.52:4.6,samples=60] (\x,{10/(\x*\x)});
  \node[ecnode,above] at (1.62,4.6) {$h_i(\cdot,u')$};
  % crossings
  \node[bundle] at (1.2,3.333) {};
  \node[bundle] at (2.5,1.6) {};
  \draw[densely dotted] (0,3.333) node[ecnode,left] {$p_i^0$} -- (1.2,3.333) -- (1.2,0);
  \draw[densely dotted] (0,1.6) node[ecnode,left] {$p_i'$} -- (2.5,1.6) -- (2.5,0);
  \node[ecnode,below,align=center,font=\scriptsize] at (1.05,0)
       {$x_i(p^0,w)$\\$h_i(p^0,u^0)$};
  \node[ecnode,below,align=center,font=\scriptsize] at (2.75,0)
       {$x_i(p',w)$\\$h_i(p',u')$};
\end{tikzpicture}
\end{center}

\topic{Step 3. Put everything together: compare EV and CV}
Both are areas under Hicksian demand, between $p_i'$ and $p_i^0$: $EV$ under
$h_i(\cdot,u')$, $CV$ under $h_i(\cdot,u^0)$.

\begin{center}
\begin{tikzpicture}[scale=0.82]
  \begin{scope}
    \fill[black!18] (0,1.6) -- plot[domain=2.5:1.7321,samples=40] (\x,{10/(\x*\x)})
                    -- (0,3.333) -- cycle;
    \ecaxes{5.4}{4.9}
    \draw[budget] plot[domain=0.95:5.0,samples=60] (\x,{4/\x});
    \draw[budget,dashed] plot[domain=1.05:2.6,samples=60] (\x,{4.8/(\x*\x)});
    \draw[budget,dashed] plot[domain=1.52:4.6,samples=60] (\x,{10/(\x*\x)});
    \draw[densely dotted] (0,3.333) node[ecnode,left] {$p_i^0$} -- (1.73,3.333);
    \draw[densely dotted] (0,1.6) node[ecnode,left] {$p_i'$} -- (2.5,1.6);
    \node[ecnode,above] at (1.62,4.6) {$h_i(\cdot,u')$};
    \node[ecnode] at (0.8,2.45) {$EV$};
  \end{scope}
  \begin{scope}[xshift=7.4cm]
    \fill[black!18] (0,1.6) -- plot[domain=1.7321:1.2,samples=40] (\x,{4.8/(\x*\x)})
                    -- (0,3.333) -- cycle;
    \ecaxes{5.4}{4.9}
    \draw[budget] plot[domain=0.95:5.0,samples=60] (\x,{4/\x});
    \draw[budget,dashed] plot[domain=1.05:2.6,samples=60] (\x,{4.8/(\x*\x)});
    \draw[budget,dashed] plot[domain=1.52:4.6,samples=60] (\x,{10/(\x*\x)});
    \draw[densely dotted] (0,3.333) node[ecnode,left] {$p_i^0$} -- (1.2,3.333);
    \draw[densely dotted] (0,1.6) node[ecnode,left] {$p_i'$} -- (1.73,1.6);
    \node[ecnode,below right] at (2.3,0.71) {$h_i(\cdot,u^0)$};
    \node[ecnode] at (0.65,2.45) {$CV$};
  \end{scope}
\end{tikzpicture}
\end{center}

Since $h_i(\cdot,u')$ lies to the right of $h_i(\cdot,u^0)$,
\[ EV>CV . \]

$EV$ and $CV$ both measure the change in indirect utility, but in different ways;
both are sound concepts.


\section{Consumer surplus}

If you used the area under Marshallian demand --- consumer surplus --- you get
neither $EV$ nor $CV$. Conceptually: what is its meaning?

$CS$ is on shaky ground. The reason is that wealth effects imply that Marshallian
demand $\neq$ Hicksian demand.

The problem disappears if there are no wealth effects. Quasilinear preferences have
this property --- no wealth effect:
\[ u(x_i,x_j)=x_j+\varphi(x_i). \]
Indifference curves move in a parallel way, and a change in $p_i$ has no wealth
effect on the consumption of $x_i$ (at an interior solution).

In this case $EV=CV=CS$, because the three curves coincide.


\begin{transcriptnotes}
  \item Notation follows Lessons 7 and 9: the original's capitals $P$, $E$, $H$,
        $X$, $V(P,w)$ are typeset as $p$, $e$, $h$, $x$, $v(p,w)$. The original's
        utility levels $V^0$, $V'$ are typeset $u^0$, $u'$, so that $v$ stays the
        indirect utility \emph{function} (as $\bar u$ in Lesson 9). $EV$, $CV$ and
        $CS$ keep their capitals. \textbf{Added:} the line defining $u^0$, $u'$ and
        the duality identities $e(p^0,u^0)=w=e(p',u')$ that the derivations use.
  \item \textbf{Clarified:} the original says ``today'' both for the old prices
        (``paid today to forgo the price change'') and for the new ones (``under
        $p'$ today''), and calls both benchmarks ``today's price''; each is replaced
        by \emph{at the old prices $p^0$} or \emph{at the new prices $p'$}.
        \textbf{Added:} the Slutsky equation and the slope comparison hold at
        $u=v(p,w)$, the comparison strictly only when $x_i>0$. ``Hicksian demand is
        flatter $\Rightarrow$ the inverse hicksian demand curve is steeper'' is
        typeset as \emph{less responsive to its own price} / \emph{steeper with
        price on the vertical axis}. Curves in the figures, written $H_i(p,V')$,
        $x_i(p,w)$, are labelled $h_i(\cdot,u')$, $x_i(\cdot,w)$: they are drawn
        against $p_i$ with the other prices held at $p_{-i}^0$.
  \item The title page reads ``Class 10'' only; the subtitle is written to match
        Lesson 9. No date is given.
  \item Page 2: ``Visiting the money metric indirect utility function to derive of
        the welfare impact'' $\to$ \emph{Revisiting \dots\ to derive the welfare
        impact}; ``price $\Delta$'' $\to$ \emph{price change}. ``the size does
        depend on $\bar p'$'' --- the prime is read as a stray mark ($\bar p$).
        ``weafare'' $\to$ \emph{welfare}; ``equiolent varation'' $\to$
        \emph{equivalent variation}.
  \item Page 3: ``the consume receives \dots\ Then he can afford.\ until $v'$''
        $\to$ \emph{then he can afford $u'$}; ``amout'' $\to$ \emph{amount}.
        In the CV text, ``If under $p'$ today $w-E(p',v^0)$'' has no verb; read as
        \emph{you give up} $w-E(p',v^0)$ (the amount is positive when $u'>u^0$).
        ``tou'' $\to$ \emph{you}. The CV definition is written with
        $\bar p-p'$ for $\bar p=p'$, and its second line ends $E(p',v^0)$, typeset
        as written.
  \item Page 4: ``EV and CV \dots\ But likely to be diff.''\ $\to$ \emph{different}.
        ``fundemental of calculee'' $\to$ \emph{fundamental theorem of calculus}.
        The integrand is written $H_i(P_1^0,P_2^0,\dots,P_i,\dots,P_i^0,V')$; the
        last price is read as $p_L^0$ (only $p_i$ varies).
  \item Page 5: ``$EV$ = are under hicksian demand curve evaluation'' $\to$
        \emph{area \dots\ evaluated at $u'$}; ``$CV$ = \dots\ $v^0$'' is a ditto
        line. The Slutsky equation's left side is written $\partial H_i(P,u)$;
        typeset $h_i(p,u)$. ``because concavity'' is completed as \emph{concavity
        of the expenditure function} (Lesson 9).
  \item \textbf{Corrected:} ``add assumption that good $i$ is a common good
        $\partial X_i(P,U)/\partial U>0$''. The property used in the Slutsky
        equation is $\partial x_i(p,w)/\partial w>0$, i.e.\ a \emph{normal} good;
        typeset that way. \textbf{Added:} the inequality
        $\partial h_i/\partial p_i>\partial x_i/\partial p_i$ is strict only when
        $x_i>0$ (it follows by subtracting the positive wealth term), and the
        parenthetical ``price on the vertical axis'' explaining why ``flatter''
        and ``steeper'' are both right.
  \item Page 7: the second duality line is written $x_i(p',E(p',u'))$; read as
        $e(p',u')$. ``Since good is an ordinary good'' $\to$ \emph{the good}.
        \textbf{Added:} the sentence stating where each Hicksian curve crosses
        Marshallian demand.
  \item Figures: the original is hand-drawn in red. Page 7 shows Marshallian
        demand $x_i(p,w)$ and two steeper curves, the right one labelled
        $H_i(p,v')$ and the left one unlabelled; the left one is labelled
        $h_i(\cdot,u^0)$ here, as the axis labels $H_i(p^0,v^0)$ require. Page 8 shows
        two panels with the area left of the Hicksian curve between the two price
        lines hatched in blue, labelled $EV$ (left) and $CV$ (right). They are
        redrawn with illustrative curves ($p=4/x$, $p=4.8/x^2$, $p=10/x^2$), grey
        fill for the hatching and dashed lines for Hicksian demand.
  \item Page 8: ``Both = area under hicksian demand'' kept;
        ``both measures change in indirect utility but in different ways sound
        concepts'' $\to$ \emph{\dots; both are sound concepts}.
  \item Page 9: ``you get either $EV$ nor $CV$'' $\to$ \emph{neither};
        ``conceputalization: what's the meanuy?''\ $\to$ \emph{Conceptually: what
        is its meaning?}; ``disapears'' $\to$ \emph{disappears}; ``this
        properties'' $\to$ \emph{this property}.
  \item Page 10: the quasilinear utility is written $V(x_i,x_j)=x_j+\hat V(x_i)$;
        typeset $u(x_i,x_j)=x_j+\varphi(x_i)$, renamed so that it does not clash
        with the indirect utility function $v$. \textbf{Corrected:} ``a $\Delta$ in $p_i$ does not impact
        consumption of $x_i$'' --- with quasilinear utility a change in $p_i$
        \emph{does} change $x_i$ (through the substitution effect); what vanishes
        is the wealth effect, so the line is typeset as ``has no wealth effect on
        the consumption of $x_i$''. The closing ``(WHEN intension)'' is not
        legible as written; it is read as \emph{at an interior solution}
        \textit{[unclear: possible reading]}.
  \item Pages of the original, for tracing back: page 1 title; page 2 the
        money-metric welfare change, the two natural benchmarks and the start of
        EV; page 3 EV and its interpretation, CV; page 4 the example, the
        Shephard's-lemma route and the EV integral; page 5 EV/CV as areas,
        step 1 (Slutsky equation, normal good); page 6 Hicksian flatter, step 2;
        page 7 duality crossings and the demand-curve graph; page 8 the EV/CV
        panels and $EV>CV$; page 9 consumer surplus and quasilinearity; page 10
        the quasilinear utility function.
\end{transcriptnotes}

\end{document}
